\section{Quadratic pseudospectrum and localized approximate states}
\label{sec:quadratic}
The quadratic pseudospectrum asks for a normalized state that minimizes
the combined residual of incompatible observables \cite{Cerjan2023}.
For a query $(E,\bm x)$, set $A_0=H-ES$, $A_j=X_j-x_jS$,
$w_0=1$, and $w_j=\kappa$ for $j=1,2,3$. The AO problem is
\begin{equation}
 Qc=\lambda Sc,\qquad
 Q=\sum_{j=0}^3 w_j^2 A_j^\dagger S^{-1}A_j,\qquad c^\dagger Sc=1 .
 \label{eq:quadratic}
\end{equation}
The minimum $\sqrt{\lambda}$ is the quadratic gap.
Its eigenvector is a joint approximate state, and several low
eigenpairs can describe a localized subspace.
The separated residuals are
\begin{equation}
 \begin{split}
 \delta E^2&=(A_0c)^\dagger S^{-1}(A_0c),\\
 \delta x_j^2&=(A_jc)^\dagger S^{-1}(A_jc),\\
 \lambda&=\delta E^2+\kappa^2\sum_j\delta x_j^2 .
 \end{split}
 \label{eq:quadratic-residuals}
\end{equation}
Expectations $c^\dagger Hc$ and $c^\dagger X_jc$ distinguish bias from
spread around the query. These states need not be eigenstates of $H$.
Within degenerate minimizing spaces, projectors or subspace observables
must be compared instead of individual coefficients.

Equation~\eqref{eq:quadratic} is positive semidefinite and covariant
under Eq.~\eqref{eq:basis-change}. Ordinary matrix products without
metric solutions would destroy that invariance.
The position residual concerns the projected position:
$X_jS^{-1}X_j$ is not an independently integrated continuum
$\langle r_j^2\rangle$. Small bases can underestimate continuum
localization errors even when the eigenproblem is solved accurately.

The dense reference uses a complete orthonormalized space.
The iterative path applies $Q$ without forming it: multiply by
$A_j$, solve with $S$, multiply by $A_j^\dagger$, and accumulate.
This avoids fill-in from explicitly squaring sparse matrices.
A complex block Rayleigh--Ritz iteration, double metric
reorthogonalization, and thick restarts address clustered low states.
The residual
\[
 \|Qc-\lambda Sc\|_{S^{-1}}
 \leq\epsilon_{\mathrm{eig}}\max(1,|\lambda|)
\]
checks an eigenpair in atomic units; it does not independently prove
that no lower state was missed.

Reusable MUMPS factors avoid forming $S^{-1}$, but the present
trial vectors and multiple-right-hand-side communication are not
fully distributed. Output includes separated residuals,
expectations, and optional complex AO coefficients.
The finite analysis covers isolated GPW/GAPW systems and post-SCF
GTH SOC. A small quadratic gap is a localization diagnostic,
not a Chern or $\mathbb Z_2$ index.

\subsection{Periodic joint localization on an AO torus}
The periodic formulation replaces the unbounded position by paired
trigonometric coordinates on a Born--von Karman torus. For cell
matrix $h$ and a full Gamma-centered $N_1\times N_2\times N_3$
property mesh, define, for each periodic lattice direction $j$,
\begin{equation}
 \begin{split}
 g_j&=\frac{2\pi}{N_j}(h^{-1})_{j,:},\qquad r_j=|g_j|^{-1},\\
 C_j&=\langle\cos(g_j\cdot r)\rangle_{\rm AO},\qquad
 D_j=\langle\sin(g_j\cdot r)\rangle_{\rm AO}.
 \end{split}
 \label{eq:quadratic-torus}
\end{equation}
These covariant matrices are assembled from the same ordered
Berry integrals and real-cell images as the periodic localizer.
For query $x$, replace each Cartesian observable in
Eq.~\eqref{eq:quadratic} by the two shifted operators
$C_j-\cos(g_j\cdot x)S$ and $D_j-\sin(g_j\cdot x)S$, both
weighted by $\kappa r_j$. With their dimensionless residuals
$\delta C_j$ and $\delta D_j$, the resulting problem satisfies
\begin{equation}
 \lambda=\delta E^2+\kappa^2\sum_{j\ \mathrm{periodic}}
 \delta\ell_j^2,\qquad
 \delta\ell_j=r_j\sqrt{\delta C_j^2+\delta D_j^2}.
 \label{eq:quadratic-chord}
\end{equation}
The length-valued chord residual is invariant under shifting the
query by a torus lattice vector. For skew cells this is a specified
reciprocal-coordinate metric, not an isotropic Cartesian variance;
changing the primitive-lattice basis can change that metric.
Finite AO projection also means that $C_j+iD_j$ need not be exactly
unitary. We do not replace it by a polar-unitarized operator.

The \texttt{FORMULATION PERIODIC} path uses the complete torus AO
space at the frozen SCF potential, independently of SCF symmetry
weights or reconstructed occupied MOs. It therefore accepts a
symmetry-reduced SCF calculation without claiming an irreducible
property eigensolver. One, two or three periodic directions are
supported; open directions are not localized by this formulation.
Output reports chord residuals and mean cosines/sines rather than
ambiguous Cartesian mean positions. Dense and DBCSR/MUMPS solvers
and the restricted post-SCF GTH-SOC construction share this tuple
of Hermitian operators. Translation-operator band unfolding is a
different joint-eigenvector problem and is not implemented here.

Small GPW/GAPW He and GTH-SOC Ne checks give serial-dense versus
two-rank-iterative gaps agreeing within $6\times10^{-15}$ hartree;
the maximum iterative eigenpair residual is
$6.7\times10^{-11}$ hartree$^2$. The He primitive-mesh and separately
converged Gamma-supercell gaps differ by less than $10^{-9}$ hartree.
Query wrapping, rigid cell rotation and equivalent SCF mesh schemes
provide additional consistency checks. Settings and limitations
are given in the Supporting Information; these are numerical
integration tests, not converged material localization predictions.
